% GENERATED by scripts/build-lecture-slides.py; do not edit. Portable profile: XeTeX, theme and fonts from the core.
\makeatletter\IfFileExists{core/beamer/beamerthememx.sty}{\def\mxcorepath{core/}\def\input@path{{core/beamer/}{core/}}}{\def\mxcorepath{../core/}\def\input@path{{../core/beamer/}{../core/}}}\makeatother
\documentclass[aspectratio=169,10pt,t]{beamer}
\usetheme{mx}
\usepackage{graphicx,booktabs}
\usepackage[french]{babel}
\setlength{\parskip}{.45em}
\hypersetup{pdfauthor={},pdfsubject={Lecture slides},pdfkeywords={},pdftitle={Équations et inéquations}}
\providecommand{\N}{\mathbb{N}}
\providecommand{\Z}{\mathbb{Z}}
\providecommand{\D}{\mathbb{D}}
\providecommand{\Q}{\mathbb{Q}}
\providecommand{\R}{\mathbb{R}}
\providecommand{\C}{\mathbb{C}}
\providecommand{\vect}[1]{\overrightarrow{#1}}
\providecommand{\vi}{\vec{\imath}}
\providecommand{\vj}{\vec{\jmath}}
\providecommand{\vk}{\vec{k}}
\providecommand{\norm}[1]{\left\lVert #1\right\rVert}
\providecommand{\vc}[2]{\begin{pmatrix}#1\\#2\end{pmatrix}}
\providecommand{\deter}[2]{\det\left(#1,#2\right)}
\providecommand{\abs}[1]{\left\lvert #1\right\rvert}
\providecommand{\dd}{\,\mathrm{d}}
\providecommand{\e}{\mathrm{e}}
\providecommand{\Card}{\operatorname{Card}}
\providecommand{\E}{\operatorname{E}}
\providecommand{\V}{\operatorname{V}}
\providecommand{\degre}{^{\circ}}
\providecommand{\un}[1]{\,\mathrm{#1}}
\providecommand{\sci}[2]{#1\times10^{#2}}
\providecommand{\pgcd}{\operatorname{PGCD}}
\providecommand{\ppcm}{\operatorname{PPCM}}
\title{Équations et inéquations}
\subtitle{Cours et exemples}
\date{}
\author{}
\begin{document}
\mxtitleframe
\begin{frame}{Objectifs}
Résoudre les équations du second degré avec paramètre, les cubiques à racine connue, les équations irrationnelles et celles avec valeur absolue.
\end{frame}
\begin{frame}{Second degré et paramètre}
Pour \(a\ne0\), \(ax^2+bx+c=0\) a un discriminant \(\Delta=b^2-4ac\). Si \(\Delta<0\), aucune racine réelle; si \(\Delta=0\), racine double \(-b/(2a)\); si \(\Delta>0\), deux racines \((-b\pm\sqrt\Delta)/(2a)\). Si \(\Delta>0\), le polynôme a le signe de \(a\) à l’extérieur des racines et le signe opposé entre elles. Si \(\Delta=0\), il a le signe de \(a\) sauf à la racine, où il est nul; si \(\Delta<0\), il a toujours le signe de \(a\).
\end{frame}
\begin{frame}{Second degré et paramètre}
La forme canonique est \(a[(x+b/(2a))^2-\Delta/(4a^2)]\). Résoudre un carré égal à \(\Delta/(4a^2)\) donne les trois cas; pour le signe, factoriser \(a(x-x_1)(x-x_2)\) dans le cas de deux racines.
\end{frame}
\begin{frame}{Second degré et paramètre}
Si \(x_1,x_2\) sont les racines, leur somme vaut \(-b/a\) et leur produit \(c/a\). Réciproquement deux nombres de somme \(S\) et de produit \(P\) sont racines de \(X^2-SX+P\).
\end{frame}
\begin{frame}{Second degré et paramètre}
Développer \(a(X-x_1)(X-x_2)\) et identifier les coefficients.
\end{frame}
\begin{frame}{Second degré et paramètre}
\((m-1)x^2-2mx+m+1=0\). Pour \(m=1\), l’équation est linéaire : \(x=1\). Pour \(m\ne1\), \(\Delta=4\) et les racines sont \(1\) et \((m+1)/(m-1)\). Ne pas appliquer la formule avant de vérifier le coefficient du carré.
\end{frame}
\begin{frame}{Degré trois et équations bicarrées}
Si \(P(r)=0\), le polynôme \(P\) est divisible par \(x-r\). Pour une cubique, le quotient est du second degré.
\end{frame}
\begin{frame}{Degré trois et équations bicarrées}
La division donne \(P(x)=(x-r)Q(x)+c\), avec reste constant. En \(x=r\), \(c=P(r)=0\).
\end{frame}
\begin{frame}{Degré trois et équations bicarrées}
Pour \(P(x)=x^3-2x^2-x+2\), on vérifie \(P(1)=0\). La division s’effectue ainsi :
\[
\begin{array}{r|l}
x^3-2x^2-x+2&x-1\\
\underline{-(x^3-x^2)}&x^2-x-2\\
-x^2-x+2&\\
\underline{-(-x^2+x)}&\\
-2x+2&\\
\underline{-(-2x+2)}&\\
0&
\end{array}
\]
Donc \(P=(x-1)(x^2-x-2)=(x-1)(x-2)(x+1)\). Une équation bicarrée se ramène au second degré en posant \(X=x^2\ge0\); chaque solution \(X>0\) fournit \(x=\pm\sqrt X\).
\end{frame}
\begin{frame}{Racines carrées : conserver les conditions}
\(\sqrt{A(x)}=B(x)\) équivaut à \(A(x)\ge0\), \(B(x)\ge0\) et \(A(x)=B(x)^2\).
Pour \(\sqrt A\le B\), il faut \(A\ge0,B\ge0,A\le B^2\).
Pour \(\sqrt A>B\), si \(B<0\), tous les points du domaine conviennent; si \(B\ge0\), comparer \(A\) à \(B^2\).
\end{frame}
\begin{frame}{Racines carrées : conserver les conditions}
Le carré est strictement croissant sur \([0;+\infty[\). Les équivalences ne valent donc qu’après avoir établi la positivité des deux membres.
\end{frame}
\begin{frame}{Racines carrées : conserver les conditions}
\(\sqrt{x+2}=x\) exige \(x\ge0\). Le carré donne \((x-2)(x+1)=0\); seule \(2\) respecte la condition. La solution \(-1\) de l’équation au carré est rejetée.
\end{frame}
\begin{frame}{Valeurs absolues et tableaux de signes}
Pour \(k\ge0\), \(|u|\le k\Leftrightarrow-k\le u\le k\), et \(|u|=k\Leftrightarrow u=k\) ou \(u=-k\). Pour une différence de valeurs absolues, séparer les intervalles où chaque expression garde un signe constant.
\end{frame}
\begin{frame}{Valeurs absolues et tableaux de signes}
Exclure les zéros du dénominateur; factoriser numérateur et dénominateur; dresser le tableau des signes; choisir les intervalles et inclure uniquement les zéros autorisés pour une inégalité large.
\end{frame}
\begin{frame}{Valeurs absolues et tableaux de signes}
Isoler d’abord les valeurs qui annulent un coefficient servant de diviseur. Résoudre ces cas dans l’équation originale; résoudre ensuite le cas général. À la fin, préciser si les solutions coïncident et si elles respectent d’éventuelles contraintes du contexte. Une formule non définie pour une valeur du paramètre ne signifie pas automatiquement « aucune solution ».
\end{frame}
\begin{frame}{Valeurs absolues et tableaux de signes}
L’équation \(mx^2-x=0\) se factorise en \(x(mx-1)=0\). Pour \(m=0\), elle devient \(-x=0\) et admet \(0\). Pour \(m\ne0\), elle admet \(0\) et \(1/m\). On ne pouvait pas écrire \(1/m\) avant d’examiner \(m=0\).
\end{frame}
\begin{frame}{Valeurs absolues et tableaux de signes}
Avec \(\sqrt{A(x)}\ge B(x)\), commencer par \(A(x)\ge0\). Si \(B(x)\le0\), l’inégalité est automatiquement vraie sur ce domaine. Si \(B(x)>0\), mettre au carré donne une comparaison équivalente. Réunir ensuite les solutions des cas; ne pas seulement lister les racines de l’équation frontière.
\end{frame}
\begin{frame}{Valeurs absolues et tableaux de signes}
Pour \(\sqrt{x+1}\ge x-1\), le domaine est \(x\ge-1\). Si \(x\le1\), le membre droit est négatif ou nul, donc l’inégalité est vraie. Si \(x>1\), le carré donne \(x+1\ge(x-1)^2\), soit \(x(x-3)\le0\); avec \(x>1\), on garde \(]1;3]\). La réunion est \([-1;3]\). Les deux bornes satisfont effectivement l’inégalité.
\end{frame}
\begin{frame}{Valeurs absolues et tableaux de signes}
Si un facteur disparaît dans un quotient, sa valeur interdite reste exclue pour l’expression initiale. Dans un tableau de signes, marquer d’abord tous les pôles, puis décider des bornes de l’ensemble de solutions.
\end{frame}
\end{document}
